% ============================================================================
% TARGET DRAFT — the paper as it should read once the causal sweep lands.
%
% READ THIS BEFORE EDITING:
%   Numbers already measured appear as plain values and are real (V0 arm,
%   diagnostics, baselines -- see evals/ and evals/diagnostics_202608/RESULTS.md).
%   Numbers NOT yet measured appear as \pending{...} and render as a loud red
%   [PENDING: ...] in the PDF. NEVER replace a \pending with a plausible-looking
%   value: fill it from a results file or leave it. A stripped placeholder is
%   exactly how a fabricated result enters a paper.
%
%   \pendingcount at the end prints how many remain.
% ============================================================================
\documentclass[runningheads]{llncs}

\usepackage{graphicx}
\usepackage{amsmath,amssymb}
\usepackage{booktabs}
\usepackage{multirow}
\usepackage{pifont}
\usepackage{xcolor}
\usepackage[hidelinks]{hyperref}

\DeclareMathOperator*{\argmax}{arg\,max}
\newcommand{\casfg}{\mathrm{CAS}_{\mathrm{fg}}}
\newcommand{\dd}{\mathrm{DD}}

% Unmeasured result. Deliberately ugly so it cannot survive to submission.
\newcounter{pend}
\newcommand{\pending}[1]{\stepcounter{pend}\textcolor{red}{\textbf{[PENDING: #1]}}}
\newcommand{\pendingcount}{\textcolor{red}{\textbf{%
  \thepend\ unresolved \texttt{\textbackslash pending} placeholders remain.}}}

\begin{document}

\title{The Route Is the Reason: Making Concept Routing\\
       Causally Decisive in 3D Medical Segmentation}
\titlerunning{The Route Is the Reason}
\author{Anonymous submission}
\authorrunning{Anonymous}
\institute{Paper ID \textcolor{red}{XXXX}}
\maketitle

\begin{abstract}
Concept bottlenecks promise intrinsic interpretability: route the prediction through a
human-readable concept layer and the concept activation \emph{is} the explanation. We show
that this promise can fail silently in dense prediction, and we show how to repair it. We
build Clinical Concept Routing (CCR), a bottleneck module that assigns every encoder token to
one of $K$ clinically-labeled experts, and train it on 3D brain-tumor segmentation. By every
standard measure it succeeds: routing correlates with the clinical concepts at $\casfg$ up to
$0.96$, and it separates necrosis, edema and enhancing tumor at AUROC $0.88$ where the best
post-hoc attribution reaches only $0.74$. Yet the routing is \emph{causally inert}. Forcing a
token through a different expert moves the prediction by $0.001$ and changes $0.2\%$ of voxel
labels; zeroing the entire bottleneck costs $1.5$ Dice. We trace this to two leaks that any
U-Net-style architecture invites---a residual path that carries features around the routing
decision, and skip connections that carry semantics around the bottleneck---and show the
experts are \emph{not} degenerate ($0.9$ pairwise divergence), so alignment metrics cannot
detect the failure. A linear probe on an ordinary segmentation encoder, trained with no
concept supervision at all, already reaches $0.81$. We then close both leaks and sweep the
resulting trade-off. \pending{one sentence: does routing become causally decisive, and at what
Dice cost} Our contributions are a diagnostic protocol that concept-bottleneck work should
adopt, evidence that the naive design fails it, and an architecture that
\pending{passes/partially passes} it. The broader lesson is methodological: \textbf{alignment
metrics do not license interpretability claims, because a perfectly aligned bottleneck can be
one the model never consults.}
\keywords{Interpretability \and Concept bottlenecks \and Causal intervention \and
Medical image segmentation \and Mixture-of-experts}
\end{abstract}

% ============================================================================
\section{Introduction}
% ============================================================================
Deep segmentation models match expert radiologists on brain-tumor benchmarks, but clinical
adoption lags because their decisions cannot be verifiably explained. The dominant remedy,
post-hoc attribution~\cite{selvaraju2017gradcam,sundararajan2017ig,lundberg2017shap}, produces
a saliency map \emph{after} the prediction and therefore approximates the mechanism rather
than reporting it. Rudin~\cite{rudin2019stop} argues the alternative: build models whose
decisions are interpretable by construction. Concept Bottleneck
Models~\cite{koh2020cbm} realize this for classification by forcing the prediction through a
human-readable concept layer, so the concept activation \emph{is} the explanation.

Extending that idea to dense prediction is attractive and, we will show, treacherous. We build
the natural extension---\textbf{Clinical Concept Routing (CCR)}: replace the encoder bottleneck
of a segmentation network with a router that assigns every token to one of $K$ clinically
labeled experts $\{$Background, Necrotic core, Edema, Enhancing tumor$\}$. The routing
distribution is a per-region, concept-labeled map, computed in one forward pass, with no
post-hoc step. It looks exactly like an intrinsic explanation should.

\paragraph{It aligns beautifully, and it does nothing.}
Trained end-to-end, CCR's routing tracks the clinical concepts closely ($\casfg$ up to $0.96$)
and discriminates them far better than any post-hoc method ($0.88$ vs.\ $0.74$). Both are the
metrics this literature reports, and by both the module is a success. But when we
\emph{intervene}---force the tokens the router assigned to necrosis through the edema expert,
changing nothing else---the prediction does not move: the target posterior shifts by $0.001$
and $0.2\%$ of voxels change label. Zeroing the bottleneck entirely costs $1.5$ Dice. The
routing is a well-trained annotation that the decoder never consults.

The failure is invisible to alignment metrics by construction. $\casfg$ measures correlation
between routing and ground truth, and correlation is compatible with an inert bottleneck.
Worse, the alignment loss trains precisely what $\casfg$ measures, so a high score confirms
the router learned its supervision, not that the network uses it.

\paragraph{Why it happens.}
We rule out the obvious explanation: the experts have \emph{not} collapsed to a shared
function---they displace tokens by $59$--$74\%$ of their norm and differ pairwise by $0.9$.
The cause is architectural, and it is two independent leaks that U-Net-style designs invite:

\begin{enumerate}
\item \textbf{Through the module.} The expert computes $h + c_k(h)$. The residual carries the
full feature vector to the decoder whichever expert fires, so routing modulates the
representation but never gates it.
\item \textbf{Around the module.} The decoder's skip connections carry encoder features
directly to the output. In our transformer backbone the CCR tap sits at stage 2 of 5, so the
two \emph{deeper}, more semantic stages bypass it entirely---the decoder can decide what tissue
it is looking at without consulting the routing at all.
\end{enumerate}

Either leak alone suffices to make routing inert, and neither is visible in $\casfg$.

\paragraph{Closing them.}
We gate both---the expert returns $c_k(h)$ alone, and a scalar $\alpha$ attenuates the
bypassing skips---and sweep $\alpha$ to trace what a genuinely causal concept bottleneck costs.
\pending{one-sentence result of the sweep}

\noindent\textbf{Contributions.}
\begin{enumerate}
\item \textbf{A diagnostic protocol for concept bottlenecks.} Three cheap tests---causal
intervention, bottleneck ablation, and a linear-probe control---that distinguish a bottleneck
the model uses from one it ignores. Alignment metrics cannot (Sec.~\ref{sec:protocol}).
\item \textbf{Evidence the naive design fails it.} A concept bottleneck with $\casfg$ up to
$0.96$ and best-in-class concept discriminability is causally inert, and we localize the cause
to two specific architectural leaks (Sec.~\ref{sec:inert}).
\item \textbf{A repair, and its price.} Gating both leaks makes the routing
\pending{causally decisive / measurably more decisive}, at a segmentation cost of
\pending{Dice delta}. We report the full trade-off curve rather than a single operating point
(Sec.~\ref{sec:repair}).
\item \textbf{A probe control that reframes the baseline.} On an encoder trained with no
concept supervision, a linear probe already separates the concepts at $0.81$---above every
post-hoc method. Concept-discriminability improvements over saliency baselines are therefore
weaker evidence than the literature treats them as (Sec.~\ref{sec:probe}).
\end{enumerate}

% ============================================================================
\section{Related Work}
% ============================================================================
\subsubsection{Post-hoc attribution.}
Grad-CAM~\cite{selvaraju2017gradcam}, integrated gradients~\cite{sundararajan2017ig},
SHAP~\cite{lundberg2017shap} and attention rollout~\cite{abnar2020rollout} all reconstruct a
saliency map after prediction. On dense prediction the extension proceeds through heuristics,
and the output is an input-\emph{sensitivity} map rather than a concept assignment. We use four
such methods as baselines and evaluate them under an identical protocol.

\subsubsection{Intrinsic interpretability and concept bottlenecks.}
CBMs~\cite{koh2020cbm} and ProtoPNet~\cite{chen2019protopnet} make the explanation part of the
computation, for scalar-output classification. Neither addresses dense prediction, where the
explanation must be per-voxel. Critically for this paper, the standard validation in this line
is \emph{concept accuracy or alignment}---does the concept layer predict the annotated
concepts---which our results show is not sufficient. A line of work on \emph{concept leakage}
in CBMs shows that soft concept representations can smuggle task information past the
bottleneck \pending{cite the concept-leakage papers -- Margeloiu et al.\ ``Do Concept
Bottleneck Models Learn as Intended?'' and Mahinpei et al.\ ``Promises and Pitfalls of
Black-Box Concept Learning Models''; VERIFY venue/year before adding to references.bib};
we report a complementary and more basic failure, in which information bypasses the bottleneck
through the \emph{architecture} rather than through the concept encoding. Positioning against
that literature matters: leakage says the concept vector carries too much, whereas we find the
concept vector is not read at all.

\subsubsection{Mixture-of-experts routing.}
Sparse MoE~\cite{shazeer2017moe,fedus2022switch,riquelme2021vmoe} introduced routing for
conditional computation. There, routing is causally decisive because the selected expert's
output \emph{is} the layer output. Our results show that transplanting the mechanism into a
skip-connected encoder--decoder silently voids that property, which is the assumption an
interpretability claim would rest on.

% ============================================================================
\section{The CCR Module}\label{sec:method}
% ============================================================================
\subsection{Router, experts, dispatch}
Let $h_t\in\mathbb{R}^D$ be the bottleneck token at spatial position $t$. A router produces
$p(t)=\mathrm{softmax}\big(W h_t/\tau + \beta\,\mathrm{sim}(h_t,\Pi)\big)\in\Delta^{K-1}$,
where $\Pi$ are $K$ learnable concept prototypes and $\tau$ is an annealed temperature. Each
concept owns an expert $E_k$. Training uses soft dispatch, inference hard dispatch:
\begin{equation}
\text{soft: } \tilde h_t=\textstyle\sum_k p_k(t)E_k(h_t),
\qquad
\text{hard: } \tilde h_t=E_{k^\ast(t)}(h_t),\; k^\ast(t)=\argmax_k p_k(t).
\label{eq:dispatch}
\end{equation}
$p(t)$ is the explanation: a per-region distribution over clinical concepts, read directly from
the forward pass.

\subsection{The two leaks}\label{sec:leaks}
Eq.~\ref{eq:dispatch} is what an intrinsic explanation requires \emph{only if} $\tilde h_t$ is
the decoder's sole source of concept information. Two standard design choices break that.

\paragraph{Residual (through).} The natural expert is a pre-norm residual MLP,
$E_k(h)=h+c_k(h)$, whose residual stabilizes early training. It also means the decoder receives
$h$ intact regardless of $k^\ast$. Routing then reweights a signal that arrives anyway. We
expose this as a switch: $E_k(h)=c_k(h)$ makes the expert transform the only path. The standard
near-zero initialization of $c_k$'s output layer must be relaxed in this mode, since $c_k$ is
now the signal rather than a correction.

\paragraph{Skips (around).} Encoder--decoder segmentation networks connect encoder stages
directly to matching decoder stages. Every such connection is a path from input to output that
does not pass the bottleneck. We attenuate them by a scalar $\alpha$, sweeping from $\alpha{=}1$
(standard) to $\alpha{=}0$ (bottleneck is the only route). Skips carry spatial detail as well as
semantics, so $\alpha$ trades boundary quality against causal control---the trade-off
Sec.~\ref{sec:repair} measures.

A related placement issue is easy to miss. In a hierarchical transformer encoder the ``bottleneck''
is ambiguous, and inserting CCR at an intermediate stage leaves \emph{deeper}, more semantic
stages flowing around it. We report this because our own first design made exactly this error.

\subsection{Objective and curriculum}
We train with $\mathcal{L}=\mathcal{L}_{\text{seg}}+\lambda_1\mathcal{L}_{\text{align}}
+\lambda_2\mathcal{L}_{\text{div}}+\lambda_3\mathcal{L}_{\text{ent}}
+\lambda_4\mathcal{L}_{\text{bnd}}$, where $\mathcal{L}_{\text{align}}$ is a foreground-only
focal BCE pushing $p_k(t)$ toward the ground-truth concept indicator,
$\mathcal{L}_{\text{div}}$ decorrelates routing vectors, $\mathcal{L}_{\text{ent}}$ is a small
entropy regularizer ($\lambda_3\le0.01$), and $\mathcal{L}_{\text{bnd}}$ is a boundary term.
A three-phase curriculum (warm-up / alignment / refinement) anneals $\tau$ from $2.0$ to $0.5$.
Note that $\mathcal{L}_{\text{div}}$ acts on routing \emph{vectors}, not on expert
\emph{functions}: nothing in the objective forces $E_a\neq E_b$. We verify functional
divergence empirically rather than assuming it (Sec.~\ref{sec:inert}).

% ============================================================================
\section{A Diagnostic Protocol for Concept Bottlenecks}\label{sec:protocol}
% ============================================================================
Alignment is necessary but not sufficient. We propose three tests, all cheap, that together
establish whether a concept bottleneck is load-bearing.

\paragraph{Definitions.}
$\casfg(k)=\mathrm{Pearson}\big(p_k(t),M_k(t)\big)$ over foreground tokens measures
routing/ground-truth \emph{alignment}. $\dd(k)=1-\mathrm{Pearson}\big(P^{\text{dec}}_k,p_k\big)$
measures how far the decoder's posterior departs from the routing.

\paragraph{Test 1: causal intervention.} Replace $k^\ast(t)$ with a chosen $b$ for the tokens
assigned to $a$, changing nothing else, and re-run the decoder. Report the change in the
posterior of class $b$ over exactly those voxels, the fraction of their labels that flip to
$b$, the change in \emph{off-target} classes (specificity), and the change \emph{outside} the
intervened region (locality). The baseline must run through the same hard-dispatch path so the
identity intervention $a\!\to\!a$ is a bit-exact no-op; then any off-diagonal effect is
attributable to re-routing alone. \emph{This test is impossible for post-hoc attribution:
there is nothing in a saliency map to intervene on.}

\paragraph{Test 2: bottleneck ablation.} Replace the module output with the raw encoder
features (\emph{bypass}) or with zeros (\emph{zero}), and measure the Dice change. If
destroying the bottleneck barely hurts, the decision does not live there and no amount of
concept alignment will make it interpretable.

\paragraph{Test 3: linear-probe control.} Fit multinomial logistic regression on the features
\emph{entering} the module and score it exactly as the router is scored. This asks whether the
concept structure is produced by the concept machinery or merely read off features that already
contained it. Run it on an ablated model trained without concept supervision for the strong
version of the control.

% ============================================================================
\section{Experiments}\label{sec:experiments}
% ============================================================================
\subsection{Setup}
BraTS 2024 glioma~\cite{menze2014brats,baid2021brats}: four co-registered MRI modalities,
labels remapped to $\{$0:Background, 1:NCR, 2:Edema, 3:ET$\}$, patient-wise
$75/12.5/12.5$ split ($n_{\text{test}}{=}168$), $z$-scored per modality, $128^3$ inputs,
$80$ epochs of AdamW ($10^{-4}$, cosine), mixed precision. We instantiate CCR on a Swin
transformer encoder with a UNETR decoder (\textbf{CCR-Net}) and on a SegResNet CNN
(\textbf{CCR-SegResNet}); the module, losses and curriculum are identical in both.

\subsection{The naive bottleneck aligns}\label{sec:aligns}
Table~\ref{tab:v0} reports the standard picture. Segmentation is competitive
($0.927/0.883/0.869$ Dice for the CNN), routing tracks the concepts closely
($\casfg$ $0.80/0.96/0.95$), and on concept-discriminability---concept $k$ versus the
\emph{other} tumor voxels---CCR reaches $0.88$ against $0.74$ for the strongest post-hoc
method, with integrated gradients and Grad-CAM near chance. Removing
$\mathcal{L}_{\text{align}}$ collapses $\casfg$ to $\approx0$ while Dice barely moves,
which we read at the time as evidence that alignment is \emph{learned}. It is; but the same
number is the first hint of the problem, since it says the segmentation does not need the
alignment.

\begin{table}[t]
\centering
\caption{The naive bottleneck by standard metrics (BraTS test, $n{=}168$). Concept-discrim.\
is concept $k$ vs.\ other tumor voxels ($0.5$=chance). By these measures the module works.}
\label{tab:v0}
\begin{tabular}{lccc}
\toprule
& Dice WT/TC/ET & $\casfg$ N/E/T & Concept-discrim. \\
\midrule
CCR-Net (Swin)        & $0.915/0.845/0.837$ & $0.69/0.94/0.93$ & $0.872$ \\
CCR-SegResNet (CNN)   & $0.927/0.883/0.869$ & $0.80/0.96/0.95$ & $0.883$ \\
\midrule
Occlusion (best post-hoc) & --- & --- & $0.741$ \\
Gradient saliency         & --- & --- & $0.601$ \\
Integrated gradients      & --- & --- & $0.583$ \\
Grad-CAM                  & --- & --- & $0.521$ \\
\bottomrule
\end{tabular}
\end{table}

\subsection{\ldots and is inert}\label{sec:inert}
\paragraph{Intervention.} Re-routing moves nothing (Table~\ref{tab:interv}). The target-class
posterior shifts by at most $0.002$ and at most $0.2\%$ of voxels change label. The effect is
not concept-specific---off-target classes move \emph{more} than the target in several cells
($0.0052$ vs.\ $0.0009$)---and not localized, since voxels outside the intervened region move
as much as those inside. The identity diagonal is exactly $0$, so this is a property of the
model and not of the instrumentation.

\begin{table}[t]
\centering
\caption{Causal intervention on CCR-Net. Re-routing the tokens assigned to the row concept
through the column concept's expert. \emph{Left:} change in the target-class posterior over the
re-routed voxels. \emph{Right:} fraction of those voxels whose predicted label became the
target. Diagonal is the identity control and is exactly zero by construction.}
\label{tab:interv}
\begin{tabular}{lccc@{\hskip 2em}ccc}
\toprule
& \multicolumn{3}{c}{$\Delta$ target posterior} & \multicolumn{3}{c}{label-flip fraction}\\
\cmidrule(lr){2-4}\cmidrule(lr){5-7}
from $\downarrow$ / to $\rightarrow$ & NCR & Edema & ET & NCR & Edema & ET \\
\midrule
NCR    & $0.0000$ & $0.0003$ & $0.0002$ & $0.0000$ & $0.0003$ & $0.0003$ \\
Edema  & $0.0009$ & $0.0000$ & $0.0012$ & $0.0008$ & $0.0000$ & $0.0007$ \\
ET     & $0.0018$ & $0.0014$ & $0.0000$ & $0.0020$ & $0.0014$ & $0.0000$ \\
\bottomrule
\end{tabular}
\end{table}

\paragraph{The experts are not degenerate.} The obvious explanation---all experts learned the
same function---is false. Experts displace tokens by $0.59$--$0.74$ of their norm and differ
pairwise by $0.84$--$1.09$ (relative $L_2$). The routing selects genuinely different
transformations; they simply do not reach the output.

\paragraph{Bottleneck ablation.} They do not reach the output because the output does not need
them. Replacing the module with raw features or with \emph{zeros} costs almost nothing:
edema Dice $0.879\to0.865$, ET $0.862\to0.806$, NCR $0.256\to0.234$. A decoder that can
segment without its bottleneck is a decoder that will ignore any concept structure placed
there.

\subsection{The probe control}\label{sec:probe}
A linear probe on the features entering the module scores $0.839$ against the router's $0.872$,
and \emph{wins} on ET. Run on the model trained with no concept supervision at all, the probe
still reaches $0.8115$---above every post-hoc method in Table~\ref{tab:v0}. The concept
structure is largely a property of a competent segmentation encoder, not of the concept
machinery. This has a consequence beyond our system: reporting that a concept module separates
concepts better than saliency baselines is weak evidence, because a logistic regression on an
unsupervised bottleneck clears that bar.

\subsection{Closing the leaks}\label{sec:repair}
We now gate both leaks and sweep $\alpha$. V0 is the naive configuration above; V1 removes the
residual; V2 and V3 additionally set $\alpha{=}0.5$ and $\alpha{=}0$.

\begin{table}[t]
\centering
\caption{The repair sweep (CCR-SegResNet). Causal control is the mean off-diagonal label-flip
fraction from Test~1; higher means the routing decision determines the output.}
\label{tab:sweep}
\begin{tabular}{llccccc}
\toprule
& residual & $\alpha$ & Dice WT/TC/ET & $\casfg$ N/E/T & Causal control & Zero-ablation $\Delta$Dice \\
\midrule
V0 & yes & $1.0$ & $0.927/0.883/0.869$ & $0.80/0.96/0.95$ & $0.002$ & $0.014$ \\
V1 & no  & $1.0$ & \pending{V1 Dice} & \pending{V1 CAS} & \pending{V1 control} & \pending{V1 abl} \\
V2 & no  & $0.5$ & \pending{V2 Dice} & \pending{V2 CAS} & \pending{V2 control} & \pending{V2 abl} \\
V3 & no  & $0.0$ & \pending{V3 Dice} & \pending{V3 CAS} & \pending{V3 control} & \pending{V3 abl} \\
\bottomrule
\end{tabular}
\end{table}

\pending{2--4 sentences: where the curve bends; whether causal control rises materially before
Dice degrades; which configuration is the recommended operating point and why}

\begin{figure}[t]
\centering
\pending{fig: causal control (x) vs Dice (y), four points V0--V3, annotated}
\caption{\textbf{The interpretability--accuracy exchange rate.} Each point is a trained model.
Moving right means the routing decision increasingly determines the prediction; moving down
means segmentation degrades. V0 is the naive concept bottleneck: maximally accurate and
causally inert.}
\label{fig:tradeoff}
\end{figure}

\paragraph{Intervention after repair.} \pending{the V3 (or best) intervention matrix, with the
same specificity and locality columns as Table~\ref{tab:interv}, showing whether re-routing now
moves the prediction toward the target concept}

\subsection{Backbone transfer}
The same module, unchanged, in a Swin transformer and a SegResNet CNN produces the same
qualitative picture: strong alignment, near-zero causal control. \pending{whether the repair
transfers to the transformer backbone as well}

% ============================================================================
\section{Discussion and Limitations}\label{sec:limitations}
% ============================================================================
\textbf{What alignment metrics can and cannot show.} $\casfg$ and concept-discriminability
measure correlation between an internal signal and annotated concepts. A module can maximize
both while contributing nothing to the prediction, and our alignment loss trains precisely what
these metrics score. Any claim that a concept layer makes a model interpretable requires an
intervention.

\textbf{No controlled accuracy comparison.} We did not train plain-backbone counterparts under
our split, so the accuracy cost of inserting the module in the V0 configuration is unquantified;
the sweep in Sec.~\ref{sec:repair} measures \emph{relative} cost across configurations, which is
the quantity the trade-off argument needs.

\textbf{Resolution.} The bottleneck is $16^3$, so the explanation is coarse and necrotic
core---tens of tokens per case---is the consistent weak spot across every metric.

\textbf{Single dataset and $K{=}4$.} Concept granularity follows the BraTS protocol; transition
zones are forced into one concept. Whether the leak analysis generalizes to other anatomies and
concept vocabularies is untested, though the two leaks are properties of the architecture family
rather than of the data.

\textbf{Scope of the explanation.} CCR explains at the level of clinical concepts---which expert
processed each region---not the router's internal reasoning. \pending{if the repair succeeds:
one sentence on what a clinician gains that the segmentation output does not already provide}

% ============================================================================
\section{Conclusion}
% ============================================================================
We set out to make the routing decision the explanation, and found that the obvious
implementation produces a concept bottleneck that is beautifully aligned and causally inert:
$\casfg$ up to $0.96$, best-in-class concept discriminability, and a prediction that does not
move when the routing is overwritten. The cause is not degenerate experts but two architectural
leaks---a residual path through the module and skip connections around it---that alignment
metrics cannot see. Closing them \pending{result}. Beyond this system, the finding is a caution
for intrinsic interpretability in dense prediction: a concept layer is only an explanation if
the model consults it, and demonstrating that requires intervening, not correlating.

\medskip\noindent\pendingcount

\bibliographystyle{splncs04}
\bibliography{references}
\end{document}
